\documentclass[11pt,a4paper]{article}
\usepackage[T1]{fontenc}
\usepackage[margin=25mm]{geometry}
\usepackage{amsmath,graphicx,booktabs,xcolor,pdfpages}
\usepackage[colorlinks=true,linkcolor=blue!50!black,hypertexnames=false]{hyperref}

% Replace these two placeholders before submitting.
\newcommand{\StudentName}{\rule{68mm}{0.45pt}}
\newcommand{\ReportDate}{\rule{68mm}{0.45pt}}


\graphicspath{{figures/}}
\setlength{\emergencystretch}{2.5em}
\setlength{\parindent}{0pt}
\setlength{\parskip}{5pt}

\title{\vspace{-12mm}
  Planning Algorithms in AI\\[2mm]
  \Large Problem Set 1 --- Discrete Planning}
\author{Gleb Kirillov\\ \small Instructor: Gonzalo Ferrer}
\date{\textbf{Date:} 14.11.2026}

\newcommand{\excerpt}[2]{%
  \smallskip\noindent
  \fcolorbox{black!35}{black!4}{%
    \begin{minipage}{\dimexpr\linewidth-2\fboxsep-2\fboxrule\relax}
    \small\textbf{#1}\par #2
    \end{minipage}}%
  \par\smallskip}
\newcommand{\quoted}[1]{\textquotedblleft#1\textquotedblright}
\newcommand{\plate}[4][\linewidth]{%
  \begin{figure}[!htbp]
    \centering
    \includegraphics[width=#1]{#2}
    \caption{#3}\label{#4}
  \end{figure}}
\newcommand{\question}[1]{\par\smallskip\textbf{#1}\par}

\newcommand{\metrics}[1]{%
  \begin{center}\small
  \begin{tabular}{lrrr}
  \toprule Search & Visited states & Share of Dijkstra & Path cost\\
  \midrule #1\bottomrule
  \end{tabular}
  \end{center}}

\begin{document}
\ifx\XeTeXversion\undefined
  \ifdefined\pdfminorversion\pdfminorversion=7\relax\fi
\else
  \special{pdf:minorversion 7}
\fi

\maketitle

\section{Task 1: Dijkstra and A* for a point robot}

\excerpt{Assignment excerpts, p.~2}{%
\quoted{Implement the Dijkstra and A* algorithms with the heuristic function
\(h(q,q_G)=0\), and plan from start to goal.}\par
\quoted{Change the heuristic function now to the \(L_1\) norm. Report the same
metrics and compare them with Dijkstra in a table.}}


\subsection{Map and search graph}

Figure~\ref{fig:point-map} shows the input occupancy grid, before search.
Let \(M(x,y)\) be the value at row \(x\) and column \(y\).
A configuration of the point robot is a two-dimensional state \(q=(x,y)\).
The graph has one vertex for each in-bounds cell with \(M(x,y)=0\).
Two vertices are joined by an undirected unit-cost edge if they differ by
one row or one column, with the other coordinate fixed. An occupied cell
has no vertex, so no edge can end in a collision. The search generates
edges by testing the four neighboring cells on demand.
The start is \(q_S=(6,6)\) and the goal is \(q_G=(85,85)\).

\plate[0.40\linewidth]{point-map.png}{Original occupancy map. Black marks
obstacles, white marks free cells, and the markers indicate the start
and goal.}{fig:point-map}

\clearpage
\subsection{Algorithms and heuristics}

The priority of a state in A* is
\begin{equation}
f(q)=g(q)+h(q,q_G),
\label{eq:astar-priority}
\end{equation}
where \(g\) is the least cost found so far from the start. Dijkstra uses
Equation~\eqref{eq:astar-priority} with the zero heuristic
\begin{equation}
h_0(q,q_G)=0.
\label{eq:zero}
\end{equation}
The informed search uses the Manhattan heuristic
\begin{equation}
h_1(q,q_G)=|x-x_G|+|y-y_G|.
\label{eq:manhattan}
\end{equation}
A unit translation can reduce Equation~\eqref{eq:manhattan} by at most
one, so it is a consistent lower bound on the remaining cost. Whenever
a cheaper path to a state is found, its queue priority and parent are
updated. The search stops when the goal is removed from the queue, then
recovers the path by following the parent pointers. This follows Lecture
L02, slides 8, 13, and 14.

A visited state is one removed from the queue for processing, including
the goal. Ties in \(f\) favor smaller \(h\) and then earlier insertion.

\subsection{Results and answers}

\metrics{Dijkstra, \(h=0\) & 6,493 & 100.0\% & 158\\
A*, \(h=L_1\) & 781 & 12.0\% & 158\\}


Figure~\ref{fig:point-search} compares the areas explored by the two searches.
\plate[0.76\linewidth]{point-search.png}{Visited cells in blue and resulting
paths in red.}{fig:point-search}

\question{How many states does A* visit compared with Dijkstra?}
A* visits 781 states, or 12.0\% of Dijkstra's 6,493. The Manhattan heuristic
in Equation~\eqref{eq:manhattan} favors cells that make progress toward the goal.
Every state processed by A* in this run had \(f=158\). The exact count
also depends on the tie-breaking rule.

\question{What is the final path cost?}
Both algorithms return 158. The path achieves the Manhattan lower bound
with 79 moves down and 79 right. The notebook saves it as
\texttt{point\_solve.gif}.

\clearpage
\section{Task 2: Configuration space of the rod}

\excerpt{Assignment excerpts, pp.~2--3}{%
\quoted{Visualize from the given data the workspace and the different rod
configurations for each discretized orientation.}\par
\quoted{Visualize the environment together with the object.}\par
\quoted{Create the C-space for the 2D environment map.}\par
\quoted{What is the size of the C-space? How do the four orientation slices
differ, and why?}}


\subsection{State variables and rod footprints}

Task 1 uses two coordinates. The rod adds an orientation index
\(k\in\{0,1,2,3\}\), so its configuration state is three-dimensional,
\(q=(x,y,k)\). The corresponding angles are \(0^\circ,45^\circ,90^\circ\),
and \(135^\circ\). Figure~\ref{fig:footprints} shows the workspace and
footprints. Figure~\ref{fig:poses} shows two valid placements.

\plate{workspace-footprints.png}{The workspace and the four supplied footprints.
The cross marks the reference cell.}{fig:footprints}
\plate[0.78\linewidth]{rod-poses.png}{Collision-free placements of the rod at
the start and goal used in Task 3.}{fig:poses}

\question{Why do four orientations suffice, and what is lost?}
A \(180^\circ\) rotation reverses the ends of the rod but leaves its
footprint unchanged. The four masks therefore exhaust the distinct
footprints on the \(45^\circ\) sampling grid. Angles between the sampled
ones are absent. The planner tests discrete poses and does not check
the area swept during a rotation.

\clearpage
\subsection{Convolution kernels and collision test}

The convolution kernels \(K_0,\ldots,K_3\) are the four binary
\(11\times11\) masks read from the supplied \texttt{rod} array, in the
same order as Figure~\ref{fig:footprints}. In each matrix, a 1 marks a cell
occupied by the rod, and row 6, column 6 is the reference cell (one-based
matrix indices). The exact kernels used in the notebook are shown below.

\begingroup\footnotesize\setlength{\arraycolsep}{2.2pt}
\[
K_0=\left[\begin{array}{@{}*{11}{c}@{}}
0 & 0 & 0 & 0 & 0 & 0 & 0 & 0 & 0 & 0 & 0 \\
0 & 0 & 0 & 0 & 0 & 0 & 0 & 0 & 0 & 0 & 0 \\
0 & 0 & 0 & 0 & 0 & 0 & 0 & 0 & 0 & 0 & 0 \\
0 & 0 & 0 & 0 & 0 & 0 & 0 & 0 & 0 & 0 & 0 \\
0 & 0 & 0 & 0 & 0 & 0 & 0 & 0 & 0 & 0 & 0 \\
1 & 1 & 1 & 1 & 1 & 1 & 1 & 1 & 1 & 1 & 1 \\
0 & 0 & 0 & 0 & 0 & 0 & 0 & 0 & 0 & 0 & 0 \\
0 & 0 & 0 & 0 & 0 & 0 & 0 & 0 & 0 & 0 & 0 \\
0 & 0 & 0 & 0 & 0 & 0 & 0 & 0 & 0 & 0 & 0 \\
0 & 0 & 0 & 0 & 0 & 0 & 0 & 0 & 0 & 0 & 0 \\
0 & 0 & 0 & 0 & 0 & 0 & 0 & 0 & 0 & 0 & 0
\end{array}\right]\qquad K_1=\left[\begin{array}{@{}*{11}{c}@{}}
0 & 0 & 0 & 0 & 0 & 0 & 0 & 0 & 0 & 0 & 0 \\
0 & 0 & 0 & 0 & 0 & 0 & 0 & 0 & 0 & 1 & 0 \\
0 & 0 & 0 & 0 & 0 & 0 & 0 & 0 & 1 & 0 & 0 \\
0 & 0 & 0 & 0 & 0 & 0 & 0 & 1 & 0 & 0 & 0 \\
0 & 0 & 0 & 0 & 0 & 0 & 1 & 0 & 0 & 0 & 0 \\
0 & 0 & 0 & 0 & 0 & 1 & 0 & 0 & 0 & 0 & 0 \\
0 & 0 & 0 & 0 & 1 & 0 & 0 & 0 & 0 & 0 & 0 \\
0 & 0 & 0 & 1 & 0 & 0 & 0 & 0 & 0 & 0 & 0 \\
0 & 0 & 1 & 0 & 0 & 0 & 0 & 0 & 0 & 0 & 0 \\
0 & 1 & 0 & 0 & 0 & 0 & 0 & 0 & 0 & 0 & 0 \\
0 & 0 & 0 & 0 & 0 & 0 & 0 & 0 & 0 & 0 & 0
\end{array}\right]
\]
\endgroup

\begingroup\footnotesize\setlength{\arraycolsep}{2.2pt}
\[
K_2=\left[\begin{array}{@{}*{11}{c}@{}}
0 & 0 & 0 & 0 & 0 & 1 & 0 & 0 & 0 & 0 & 0 \\
0 & 0 & 0 & 0 & 0 & 1 & 0 & 0 & 0 & 0 & 0 \\
0 & 0 & 0 & 0 & 0 & 1 & 0 & 0 & 0 & 0 & 0 \\
0 & 0 & 0 & 0 & 0 & 1 & 0 & 0 & 0 & 0 & 0 \\
0 & 0 & 0 & 0 & 0 & 1 & 0 & 0 & 0 & 0 & 0 \\
0 & 0 & 0 & 0 & 0 & 1 & 0 & 0 & 0 & 0 & 0 \\
0 & 0 & 0 & 0 & 0 & 1 & 0 & 0 & 0 & 0 & 0 \\
0 & 0 & 0 & 0 & 0 & 1 & 0 & 0 & 0 & 0 & 0 \\
0 & 0 & 0 & 0 & 0 & 1 & 0 & 0 & 0 & 0 & 0 \\
0 & 0 & 0 & 0 & 0 & 1 & 0 & 0 & 0 & 0 & 0 \\
0 & 0 & 0 & 0 & 0 & 1 & 0 & 0 & 0 & 0 & 0
\end{array}\right]\qquad K_3=\left[\begin{array}{@{}*{11}{c}@{}}
0 & 0 & 0 & 0 & 0 & 0 & 0 & 0 & 0 & 0 & 0 \\
0 & 1 & 0 & 0 & 0 & 0 & 0 & 0 & 0 & 0 & 0 \\
0 & 0 & 1 & 0 & 0 & 0 & 0 & 0 & 0 & 0 & 0 \\
0 & 0 & 0 & 1 & 0 & 0 & 0 & 0 & 0 & 0 & 0 \\
0 & 0 & 0 & 0 & 1 & 0 & 0 & 0 & 0 & 0 & 0 \\
0 & 0 & 0 & 0 & 0 & 1 & 0 & 0 & 0 & 0 & 0 \\
0 & 0 & 0 & 0 & 0 & 0 & 1 & 0 & 0 & 0 & 0 \\
0 & 0 & 0 & 0 & 0 & 0 & 0 & 1 & 0 & 0 & 0 \\
0 & 0 & 0 & 0 & 0 & 0 & 0 & 0 & 1 & 0 & 0 \\
0 & 0 & 0 & 0 & 0 & 0 & 0 & 0 & 0 & 1 & 0 \\
0 & 0 & 0 & 0 & 0 & 0 & 0 & 0 & 0 & 0 & 0
\end{array}\right]
\]
\endgroup

Index the entries of each kernel by offsets \(u,v\in\{-5,\ldots,5\}\).
Extend \(M\) to \(M_{\mathrm{pad}}\) by setting every out-of-bounds cell
to 1. For each map center, compute
\begin{equation}
S_k(x,y)=(M_{\mathrm{pad}} * K_k)(x,y)
  =\sum_{u=-5}^{5}\sum_{v=-5}^{5}
  M_{\mathrm{pad}}(x-u,y-v)\,K_k(u,v).
\label{eq:convolution}
\end{equation}
Convolution reverses \(u,v\), but the supplied kernels satisfy
\(K_k(u,v)=K_k(-u,-v)\). Equation~\eqref{eq:convolution} therefore counts
exactly the occupied rod cells that touch an obstacle or lie outside the
grid. The notebook evaluates it using \texttt{scipy.signal.convolve2d}
with \texttt{mode='same'}, \texttt{boundary='fill'}, and \texttt{fillvalue=1}.
The output slice keeps the map size. It marks \((x,y,k)\) blocked when
\(S_k(x,y)>0\) and free when \(S_k(x,y)=0\). The notebook independently
checks the result by shifting and combining obstacle masks at the occupied
footprint offsets.

\clearpage
\subsection{Orientation slices and graph edges}

\plate{configuration-slices.png}{Four configuration-space slices. Black marks
workspace obstacles. Orange marks additional cells blocked by the rod footprint.
White marks free configurations. Each panel fixes one orientation index.}{fig:slices}

\begin{center}\small
\begin{tabular}{lrr}
\toprule Orientation & Free configurations & Connected components\\
\midrule
\(0^\circ\) & 4,677 & 2\\
\(45^\circ\) & 3,926 & 4\\
\(90^\circ\) & 4,574 & 1\\
\(135^\circ\) & 3,909 & 5\\
\midrule Total & 17,086 & \\
\bottomrule
\end{tabular}
\end{center}


\question{What is the size of the configuration space?}
There are \(100\times100\times4=40{,}000\) candidate configurations, of
which 17,086 (42.7\%) are free. Each panel of Figure~\ref{fig:slices}
is a two-dimensional cut through this three-dimensional state space,
with \(k\) held fixed.

\question{How do the slices differ, and why?}
A horizontal footprint extends obstacles along columns and a vertical
one along rows. Diagonal footprints extend them along both coordinates,
thickening all four border strips. The two diagonal slices have the
fewest free configurations. For example, the passage near column 66 is
open at \(90^\circ\) and blocked at \(0^\circ\). The table counts components
using only the four translations within each individual slice.

\question{How are the slices connected?}
Each free \((x,y,k)\) is a vertex. A translation joins neighboring
free cells within the same slice. An in-place turn joins free cells at
the same \((x,y)\) in adjacent slices. Specifically, slice 1 connects to
2 and 4, slice 2 to 1 and 3, slice 3 to 2 and 4, and slice 4 to 3 and 1.
These are one-based slice numbers. A connection exists only when both
endpoint configurations are free. Each edge has unit cost and can be
traversed in either direction. This gives one graph with 17,086 vertices,
more than the \(6{,}859\) free-cell vertices in the point-robot graph.

\clearpage
\section{Task 3: A* with position and orientation}

\excerpt{Assignment excerpts, p.~3}{%
\quoted{On this first iteration, use the heuristic function \(h(q,q_G)=0\),
which is equivalent to the Dijkstra algorithm.}\par
\quoted{Change the heuristic function now to be \(h(q,q_G)=L_1\) norm of the
\(x,y\) components.}\par
\quoted{Propose a heuristic function \(h(q,q_G)\) that includes orientation.
Is it admissible? Justify your answer.}}


\subsection{Heuristics on the rod graph}

We search the graph from Task 2 from \((6,6,2)\) to \((85,85,0)\).
Dijkstra uses Equation~\eqref{eq:zero} and the position-only A* search
uses Equation~\eqref{eq:manhattan}, ignoring orientation. The third
search uses
\begin{equation}
h_{\mathrm{rod}}(q,q_G)=h_1(q,q_G)+d_\theta(k,k_G),
\label{eq:rod-heuristic}
\end{equation}
where
\begin{equation}
d_\theta(k,k_G)=\min\bigl(\lvert k-k_G\rvert,\;4-\lvert k-k_G\rvert\bigr).
\label{eq:four-turns}
\end{equation}
The orientation indices form the cycle \(0\to1\to2\to3\to0\).
There are two arcs from \(k\) to a different \(k_G\). The arc through
consecutive indices without wraparound has \(\lvert k-k_G\rvert\) edges.
The other arc uses the remaining edges of the cycle, so it has
\(4-\lvert k-k_G\rvert\) edges. Equation~\eqref{eq:four-turns} takes the
shorter route because the robot can turn in either direction. For
example, \(3\to0\) takes one turn even though the indices differ by three.
When \(k=k_G\), the formula gives zero.

\question{Is the orientation heuristic admissible?}
Yes. Ignoring obstacles, every solution needs at least \(h_1\) translations
and at least \(d_\theta\) turns. A single action can contribute to only
one of those counts. Their sum in Equation~\eqref{eq:rod-heuristic} is
therefore a lower bound. Each graph edge changes only one term by at most
one, so the sum is also consistent.

\clearpage
\subsection{Results and answers}

\metrics{Dijkstra, \(h=0\) & 16,323 & 100.0\% & 164\\
A*, \(h=L_1\) & 7,720 & 47.3\% & 164\\
A*, \(h=L_1+d_\theta\) & 4,572 & 28.0\% & 164\\}


\plate{rod-search.png}{Search traces projected onto the workspace. Blue shades
count visited orientations, red lines show paths, and purple diamonds mark
turns.}{fig:rod-search}

\question{How many states are visited, and what is the final cost?}
Dijkstra visits 16,323 states. A* with Manhattan distance visits 7,720
(47.3\% of Dijkstra's count). Including rotation reduces this to
4,572 (28.0\%). All three return the same optimal cost, 164.
Figure~\ref{fig:rod-search} shows the difference in the explored regions.
The notebook saves the Dijkstra plan as \texttt{rod\_solve.gif}.

\question{How do these results relate to Task 1?}
The rod path uses the same minimum of 158 translations and adds six turns,
giving a total cost of 164. Equation~\eqref{eq:manhattan} is exact at the
point robot's start. For the rod, Equation~\eqref{eq:rod-heuristic}
raises the initial lower bound from 158 to 160, below the optimal cost of 164.

\clearpage
\section{Task 4: Forward-only motion (additional experiment)}

\textit{Student-designed extension.} I designed this experiment to show
how a coarse angular discretization can make a planning problem
unsolvable and to examine the trade-off between finding a feasible route
and keeping the configuration space small. The rod can move only
forward along its current heading, but it can still rotate in place.
I compare the four geometric orientations supplied in PS1 with a finer
angular discretization, measuring both whether the goal is reachable
and how much the configuration space grows.

\excerpt{Extension question}{%
\quoted{Can the rod reach the goal if it can only move forward along its current
orientation, while retaining in-place rotations? Does a finer angular
discretization admit a route?}}


\subsection{Search with the supplied angular discretization}

The four geometric orientations give eight directed headings, since
the two ends of each footprint define opposite forward directions.
A state can move one step forward along its heading or turn \(45^\circ\)
in place. Each action costs one, including a diagonal step. The collision
test uses the destination footprint. A forward edge is directed because
the reverse step would require facing the other way.

Dijkstra still uses Equation~\eqref{eq:zero}. For the two A* runs we use
\begin{equation}
h_\infty(q,q_G)=\max\{\lvert x-x_G\rvert,\lvert y-y_G\rvert\}.
\label{eq:chebyshev}
\end{equation}
A diagonal step can reduce both coordinate differences by one, so
Equation~\eqref{eq:chebyshev} bounds the number of forward steps. The
second heuristic adds the shortest turn count on the eight-heading cycle:
\begin{equation}
h_{\mathrm{forward}}(q,q_G)=
h_\infty(q,q_G)+\min\bigl(\lvert k-k_G\rvert,\;8-\lvert k-k_G\rvert\bigr).
\label{eq:forward-heuristic}
\end{equation}
The same two-arc argument as in Equation~\eqref{eq:four-turns} applies.
Each action changes at most one term by at most one, so both
Equations~\eqref{eq:chebyshev} and~\eqref{eq:forward-heuristic} are consistent.

\subsection{Results and answers}

\plate[0.70\linewidth]{reachable-headings.png}{Reachable centers on the full map
and reachable headings near the start.}{fig:reachable}

\clearpage
\subsection{Initial heading and the blockage}

\question{Does the original start reach the goal, and where does it get stuck?}
Dijkstra and both A* runs exhaust the same 14 configurations in six cells
from the \(90^\circ\) start. The first forward step hits the upper wall.
Turning toward \(45^\circ\) is legal. Turning toward \(135^\circ\) intersects
a nearby obstacle. Figure~\ref{fig:reachable} shows the reachable region.

\question{Would a different initial heading make the goal reachable?}
\begin{center}\small
\begin{tabular}{lrrl}
\toprule Initial headings at \((6,6)\) & Reachable states & Centers & Goal center\\
\midrule
\(0^\circ,45^\circ,90^\circ\) & 14 & 6 & Unreachable\\
\(180^\circ,225^\circ,270^\circ\) & 5 & 2 & Unreachable\\
\(135^\circ,315^\circ\) & -- & -- & Initial pose blocked\\
\bottomrule
\end{tabular}
\end{center}

Obstacles can block in-place turns, so the initial heading changes the
reachable set. None of the six valid initial headings reaches the goal center.

\plate[0.9\linewidth]{blocked-frames.png}{From left to right, the rod starts at \(90^\circ\),
hits the upper wall when trying to move forward, and overlaps an obstacle
when trying to turn to \(135^\circ\). Red marks the overlap.}{fig:blocked}

\clearpage
\subsection{Finer headings and the recovered route}

The refined model uses 16 primitive integer directions with components of
magnitude at most two, including slopes \(1:2\) and \(2:1\). It preserves
the original masks and rasterizes new footprints as centered line segments.
For longer forward moves, it checks the destination and grid cells
surrounding the midpoint. The result applies to this discrete collision model.

\begin{center}\small
\begin{tabular}{lrr}
\toprule Heading model & Visited states & Path cost\\
\midrule
8 headings & 14 & Unreachable\\
16 headings, original forward directions & 27 & Unreachable\\
16 headings, all forward directions & 52,321 & 88\\
\bottomrule
\end{tabular}
\end{center}

\question{Does finer discretization recover a route, and what change makes it possible?}
Dijkstra finds a route with 72 forward steps and 16 turns. Its first
translation goes from \((6,6)\) to \((5,8)\) at about \(26.6^\circ\). With
the same 16 headings but only the original forward directions, no route
exists. The successful refinement adds both turning angles and forward
directions. Figure~\ref{fig:blocked} shows the original blockage, and
Figure~\ref{fig:fine} shows the recovered route. The result isolates
the effect of enabling forward motion at the new angles within this
discrete model. Path cost counts actions within the respective model.

\plate[0.83\linewidth]{fine-frames.png}{The rod leaves the start along a new heading, crosses
the workspace, and reaches the goal.}{fig:fine}

\clearpage
\subsection{Conclusions}

\question{How much does the configuration space grow?}
The state still has three coordinates \((x,y,k)\). The refinement
increases the number of values of \(k\) while leaving the spatial grid
unchanged.

\begin{center}\small
\begin{tabular}{lrrr}
\toprule Model & Angular values & Candidate states & Free states\\
\midrule
Task 3, undirected footprint orientations & 4 & 40,000 & 17,086\\
Task 4, coarse directed headings & 8 & 80,000 & 34,172\\
Task 4, refined directed headings & 16 & 160,000 & 64,258\\
\bottomrule
\end{tabular}
\end{center}

From 8 to 16 headings, the candidate space doubles
(\(100\times100\times8\) to \(100\times100\times16\)). The number of free
states, which are the vertices of the search graph, increases by
30,086 (88.0\%). Relative to the four-orientation model in Task~3,
the refined graph has four times as many candidate states.

\question{What trade-off does the refinement show?}
The four supplied footprint angles, represented by eight directed
headings, make this forward-only instance unsolvable. Sampling the
angle more finely and allowing forward steps at the new headings
restores a route of 88 actions. In this model, the price of planning
with these extra directions is a larger configuration space.

\section*{Code availability}
The notebook, report source, and animations are maintained in the
\href{https://github.com/glebuben/planning-2026-T2}{planning-2026-T2 GitHub repository}
on branch \texttt{ps1}, under \texttt{PS/PS1/solution}.

\section*{AI usage acknowledgment}
A large language model (LLM) was used to generate visualizations, assist
with debugging the code, and improve the grammar and wording of the conclusions.

\clearpage
\phantomsection
\pdfbookmark[0]{Appendix A. Original problem set}{assignment}
\includepdf[pages=-,pagecommand={\thispagestyle{empty}}]{../../ps1.pdf}

\end{document}
